\begin{center}\begin{tikzpicture}[font=\small,>=stealth,line join=round]
\draw[gray!40,rounded corners=3pt] (0,0) rectangle (14.2,8.3);\node at (7.1,7.700000000000001) {例10：定义域中间有缺口，左右仍成对};
\begin{scope}[shift={(4.1,2.307142857142857)},x=1.0714285714285714cm,y=1.3142857142857143cm]\draw[->,gray] (-2.8,0)--(2.8,0) node[right]{$x$};\draw[->,gray] (0,-0.5)--(0,3) node[above]{$y$};\draw[blue,thick,solid,domain=-2.5:-1,samples=180] plot(\x,{sqrt((\x)^2-1)});\draw[blue,thick,solid,domain=1:2.5,samples=180] plot(\x,{sqrt((\x)^2-1)});\fill[orange!85!black] (-1,0) circle(2.2pt);\fill[orange!85!black] (1,0) circle(2.2pt);\end{scope}\node[anchor=west,align=left,text width=5.0cm] at (8.4,4.0) {$F=\sqrt{x^2-1}$\\[10pt]$D=(-\infty,-1]\cup[1,+\infty)$\\[10pt]两支左右对应\\[10pt]整体仍是偶函数};
\node[text width=12.6cm,align=center] at (7.1,.65) {端点实心，因为根号允许被开方数为0；缺口本身不破坏定义域对称。};
\end{tikzpicture}\end{center}
